\begin{abstract}
\label{chap:概要}\par

近年の Vision-Language-Action（VLA）モデルは移動タスクの達成性能を向上させている．一方で，「どのように移動するか」というスタイル制御は十分に扱われていない．本研究では，日本語オノマトペを用いてロボット歩行の質感を指示する手法を提案する．提案手法は，対照学習と強化学習からなる 2 段階構成である．Phase 1 では，HOYO データセットの 2D 歩容系列とオノマトペを用い，MotionCLIP を基盤とした共有潜在空間を学習する．このとき，VAE による再構成損失と教師あり対照損失を併用し，言語埋め込みと動作埋め込みの整合を獲得する．Phase 2 では，Isaac Lab 上の H1 ヒューマノイドに対して PPO で歩行方策を学習し，ロボット動作の潜在表現と教師潜在のコサイン類似度をスタイル報酬として与える．

実験では，対照学習の主指標として検索再現率指標 $\mathrm{avg\_r@1}$ を用い，$\mathrm{avg\_r@1}=0.356\pm0.224$ を得た．また，未知語の埋め込みは意味的に近い既知ラベル近傍に配置され，言語-動作対応の汎化傾向を確認した．強化学習では，速度指令値 $0.5\,\mathrm{m/s}$ をスタイル報酬のみでその場足踏みに収束することを避ける補助条件として用い，主要評価をオノマトペごとの相対速度差と見た目の妥当性に置いた．その結果，実速度は $0.148$--$0.317\,\mathrm{m/s}$ に分布し，「すたすた」は相対的に高速，「どっしどっし」「のろのろ」は相対的に低速となる傾向を確認した．以上より，オノマトペの印象と整合する歩容スタイル付与の有効性を示した．

\end{abstract}
